\documentclass[ ../main.tex]{subfiles}
\providecommand{\mainx}{..}
\begin{document}
\section{The singular hash map}
\label{sec:shm}

We present the \emph{singular hash map}, a concrete construction that realizes the entropy map framework from \cref{sec:entropy-maps} and achieves the information-theoretic space lower bound asymptotically.
The construction uses universal hashing~\cite{carterWegman1979} and is related to perfect hashing techniques~\cite{fredmanKomlosSzemeredi1984} and succinct retrieval data structures~\cite{dietzfelbinger2008}.

\subsection{Construction algorithm}
\label{sec:shm-construction}

\begin{algorithm}
	\caption{Implementation of \protect\MakeSingularHashMap over a universal set $\Set{X}$}
	\label{alg:makeset}
	\DontPrintSemicolon
	\SetKwProg{func}{function}{}{}
	\KwIn
	{
		$\Set{A}$ is a subset of a universal set $\Set{X}$.
		$\fprate$ is the false positive rate.
		$t$ is the maximum number of trials.
	}
	\KwOut
	{
		$\ASet{A}[\fprate]$, a rate-distorted approximate set of $\Set{A}$.
	}
	\func{\MakeSingularHashMap{$\Set{A}$, $\fprate$, $t$}}
	{
		$k \coloneqq \lceil -\log_2 \fprate \rceil$\;
		$\Set{A}_{\BitSet^*} \coloneqq \SetBuilder{\Encode_{\Set{X}}(a) \in \BitSet^*}{a \in \Set{A}}$\;
		\For{$n \gets 0$ \KwTo $t$}
		{
			\For{$j \gets 1$ \KwTo $2^n$}
			{
				$\found \gets \True$\;
				$b_n \gets $ a bit string randomly drawn from $\BitSet^n$ without replacement\;
				$h_k \gets \Nothing$\;
				\For{$x \in \Set{A}_{\BitSet^*}$}
				{
					\If{$h_k = \Nothing$}
					{
						$h_k \gets \trunc(\hash(x \cat b_n), k)$\;
					}
					\ElseIf{$h_k \neq \trunc(\hash(x \cat b_n), k)$}
					{
						$\found \gets \False$\;
					}
				}
				\If{\found}
				{
					\Return $\Tuple{h_k, b_n}$\;
				}
			}
		}
	}
\end{algorithm}

The space required for the \emph{singular hash map} is of the order of the length $n$ of the bit string $b$. Therefore, for space efficiency, the algorithm exhaustively searches for a bit string in the order of increasing size $n$.

\subsection{Probability of success}
\label{sec:shm-prob}

The \emph{singular hash set} proves that the keys in $\APFun{f}$ are an \emph{approximate set} of the keys in $\Fun{f}$.

\begin{theorem}
\label{thm:shm-success}
Let $\fprate = 2^{-k}$ for some positive integer $k$.
The probability that a particular bit string $b \in \BitSet^*$ successfully codes the \emph{singular hash map} is given by
\begin{equation}
    p(m,\fprate,\mu) = \frac{\fprate^{m-1}}{2^{m \mu}}\,.
\end{equation}
\end{theorem}
\begin{remark}
The restriction to $\fprate = 2^{-k}$ ensures that the hash truncation $\trunc(\cdot, k)$ yields an exact false positive rate.
For arbitrary $\fprate \in (0,1)$, one chooses $k = \lceil -\log_2 \fprate \rceil$ and obtains $2^{-k} \leq \fprate$; the result then generalizes via standard approximation arguments.
\end{remark}
\begin{proof}
For a particular bit string $b_n$, the joint probability that every key in $\Fun{f}$ collides and for each key $k$ in $\Fun{f}$, the concatenation of $1$ and $k$ has a hash in which the first $\BL(v)$ bits collide with mapped value $v$ is given by the following reasoning.

Suppose we have a set $\Fun{f} \coloneqq \{(k_1,v_1),\ldots,(k_m,v_m)\}$ and $k_1$ hashes to $y = \hash(k_1) \mod t$, where $\hash \colon \BitSet^* \to \BitSet^N$ is a random hash function that uniformly distributes over its domain of $2^N$ possibilities and $N < \max{\BL(v_1),\ldots,\BL(v_m)}$.

The probability that $\hash(1 \cat k_1) \mod \BL(v_1) = v_1$ is given by
\begin{equation}
    \frac{1}{2^{\BL(v_1)}}\,.
\end{equation}
Since $y$ is a particular element in $\BitSet^k$, the probability that $k_j, j \neq 1$, hashes to $y$ is given by
\begin{equation}
    \frac{1}{2^k}\,.
\end{equation}
The probability that $1 \cat k_j, j \neq 1$ hashes to $v_j$ is given by
\begin{equation}
    \frac{1}{2^{\BL(v_j)}}\,.
\end{equation}
Since $\hash$ is a random hash function, every one of these probabilities are independent, and thus the joint probability is just the product given by
\begin{align}
    p &= \frac{1}{2^{k(m-1)}} \prod_{j=1}^{m} \frac{1}{2^{\BL(v_j)}}\\
      &= \fprate^{m-1} \frac{1}{2^{\sum_{j=1}^{m} \BL(v_j)}}\,.
\end{align}
The average bit length per key is given by
\begin{equation}
    \mu = \frac{1}{m} \sum_{j=1}^{m} \BL(v_j)
\end{equation}
and thus the probability $p$ may be parameterized as
\begin{equation}
    p = \frac{\fprate^{m-1}}{2^{m \mu}}\,.
\end{equation}
\end{proof}

\subsection{Space complexity}
\label{sec:shm-space}

\begin{theorem}
\label{thm:shm-space}
Given \emph{false positive} and \emph{false negative} domain error rates given by $\fprate$ and $\fnrate$ respectively, and a mean bit length $\mu$ of encodings of elements in the image of $\APFun{f}$, the \emph{expected} bit length of the Singular Hash Map asymptotically obtains a space complexity with a lower-bound given by
\begin{equation}
    -(1 - \fnrate) \log_2 \fprate + (1 - \fnrate) \mu \; \si{bits \per element}\,,
\end{equation}
which occurs when the encodings have uniformly distributed bit lengths, and an upper-bound given by
\begin{equation}
    -(1 - \fnrate) \log_2 \fprate + \mu \; \si{bits \per element}\,,
\end{equation}
which occurs when the encodings have degenerate bit lengths. The greater the entropy of the bit lengths, the smaller the expected bit length given a mean of $\mu$.
\end{theorem}
\begin{proof}
The space required for the singular hash set found by \Cref{alg:makeset} is of the order of the length $n$ of the bit string $b_n$.
Therefore, for space efficiency, the algorithm exhaustively searches for a bit string in the order of increasing size $n$.

We are interested in the first case when a success occurs, which is a geometric distribution with probability of success $p$ as given by
\begin{equation}
    \RV{Q} \sim \geodist\!\left(p = \frac{\fprate^{m-1}}{2^{m \mu}}\right)\,.
\end{equation}
The expected number of trials for the geometric distribution is given by
\begin{equation}
\label{eq:exp_trials}
    \Expect{\RV{Q}} = \frac{2^{m \mu}}{\fprate^{m-1}}\,.
\end{equation}
The $n$-th trial uniquely maps to a bit string of length $\ell = \lfloor \log_2 n \rfloor$.
By Jensen's inequality (since $\log_2$ is concave),
$\Expect{\log_2 \RV{Q}} \leq \log_2 \Expect{\RV{Q}}$, so
\begin{equation}
    \Expect{\log_2 \RV{Q}} \leq \log_2 \frac{2^{m \mu}}{\fprate^{m-1}}\; \si{bits}\,.
\end{equation}
The approximation error is $O(1/m)$ per element and vanishes as
$m \to \infty$, justifying the asymptotic claim below.
Dividing by $m$ and simplifying results in
\begin{equation}
    -\frac{m-1}{m} \log_2 \fprate + \mu\; \si{bits \per element}\,.
\end{equation}
Asymptotically, as $m \to \infty$, the bits per element converges to
\begin{equation}
\label{eq:shm-zero-fnr}
    -\log_2 \fprate + \mu\,.
\end{equation}

Now suppose a fraction $\fnrate$ of the $m$ keys are dropped (false negatives), so that only $(1-\fnrate)m$ keys are successfully stored.
The key-discrimination cost applies only to the $(1-\fnrate)m$ stored elements, contributing $-(1-\fnrate)\log_2 \fprate$ bits per element of the original map.
For the value-encoding cost, we consider two cases.
In the lower-bound case (uniformly distributed bit lengths), the value bits scale proportionally with the number of stored elements, contributing $(1-\fnrate)\mu$ bits per element.
In the upper-bound case (degenerate bit lengths), the encoding overhead for each stored value remains $\mu$ bits regardless of how many elements are stored, because the bit length is constant and cannot be amortized; this contributes $\mu$ bits per element.
Combining these gives the bounds in the theorem statement.\end{proof}

\begin{corollary}
When $\fnrate = 0$ and $\mu = 0$ (the case of an approximate set with no values to store), the space complexity reduces to $-\log_2 \fprate$ bits per element, which is the information-theoretic lower-bound for first-order random positive approximate sets over countably infinite universes~\cite{cover2006}.
\end{corollary}

\end{document}
